\section*{Capítulo \ref{ch:g9:inside-the-atom} --- Dentro do átomo}

\begin{solution}{exo:g9:inside-the-atom:1}
No \omterm{def:g9:inside-the-atom:nucleus}{núcleo}: \omterm{def:g9:inside-the-atom:proton}{prótons} (positivos) e \omterm{def:g9:inside-the-atom:proton}{nêutrons} (sem carga). Em volta dele:
\omterm{def:g9:inside-the-atom:nucleus}{elétrons} (negativos).
\end{solution}

\begin{solution}{exo:g9:inside-the-atom:2}
\ce{^{16}_{8}O}: 8 \omterm{def:g9:inside-the-atom:proton}{prótons}, 8 \omterm{def:g9:inside-the-atom:proton}{nêutrons}, 8 \omterm{def:g9:inside-the-atom:nucleus}{elétrons}.
\ce{^{27}_{13}Al}: 13, 14, 13. \ce{^{40}_{20}Ca}: 20, 20, 20.
\end{solution}

\begin{solution}{exo:g9:inside-the-atom:3}
Ele tem tantos \omterm{def:g9:inside-the-atom:nucleus}{elétrons} (carga $-e$ cada um) quanto \omterm{def:g9:inside-the-atom:proton}{prótons} (carga $+e$
cada um): as cargas se anulam.
\end{solution}

\begin{solution}{exo:g9:inside-the-atom:4}
$A = 26 + 30 = 56$: \ce{^{56}_{26}Fe}.
\end{solution}

\begin{solution}{exo:g9:inside-the-atom:5}
Sim: os dois têm $Z = 17$, logo os dois são cloro. Eles têm \omterm{def:g9:inside-the-atom:atomic-number}{números de
massa} diferentes (18 e 20 \omterm{def:g9:inside-the-atom:proton}{nêutrons}): são \omterm{def:g9:inside-the-atom:isotope}{isótopos}.
\end{solution}

\begin{solution}{exo:g9:inside-the-atom:6}
Não: eles têm \omterm{def:g9:inside-the-atom:atomic-number}{números atômicos} diferentes (6 e 7), logo são elementos
diferentes, o carbono e o nitrogênio. Os \omterm{def:g9:inside-the-atom:isotope}{isótopos} têm o mesmo $Z$.
\end{solution}

\begin{solution}{exo:g9:inside-the-atom:7}
$56 \times 1.67 \times 10^{-27} \approx \qty{9.35e-26}{kg}$.
\end{solution}

\begin{solution}{exo:g9:inside-the-atom:8}
Um \omterm{def:g7:atoms-and-molecules:atom}{átomo} é principalmente espaço vazio: o \omterm{def:g9:inside-the-atom:nucleus}{núcleo} é minúsculo, por isso a
maioria das partículas não encontra nenhum \omterm{def:g9:inside-the-atom:nucleus}{núcleo} no caminho e atravessa
em linha reta.
\end{solution}

\begin{solution}{exo:g9:inside-the-atom:9}
$10^{-10} / 10^{-15} = 10^{5}$: cem mil vezes menor.
\end{solution}

\begin{solution}{exo:g9:inside-the-atom:10}
A química depende dos \omterm{def:g9:inside-the-atom:nucleus}{elétrons}, e os \omterm{def:g9:inside-the-atom:isotope}{isótopos} de um elemento têm o mesmo
número de \omterm{def:g9:inside-the-atom:nucleus}{elétrons} (o mesmo $Z$).
\end{solution}

\begin{solution}{exo:g9:inside-the-atom:11}
$26 \times 9.11 \times 10^{-31} = \qty{2.37e-29}{kg}$. Fração:
$2.37 \times 10^{-29} / 9.35 \times 10^{-26} \approx 2.5 \times
10^{-4}$, cerca de \qty{0.025}{\percent} da massa.
\end{solution}

\begin{solution}{exo:g9:inside-the-atom:12}
6915 \omterm{def:g7:atoms-and-molecules:atom}{átomos} de cobre-63 e 3085 \omterm{def:g7:atoms-and-molecules:atom}{átomos} de cobre-65. Média:
$(6915 \times 63 + 3085 \times 65)/10\,000 = 63.6$ \omterm{def:g9:inside-the-atom:proton}{núcleons}.
\end{solution}

\begin{solution}{pb:g9:inside-the-atom:1}
\textbf{1.} \ce{^{35}_{17}Cl} e \ce{^{37}_{17}Cl}.

\textbf{2.} Cloro-35: 17 \omterm{def:g9:inside-the-atom:proton}{prótons}, 18 \omterm{def:g9:inside-the-atom:proton}{nêutrons}, 17 \omterm{def:g9:inside-the-atom:nucleus}{elétrons}. Cloro-37:
17 \omterm{def:g9:inside-the-atom:proton}{prótons}, 20 \omterm{def:g9:inside-the-atom:proton}{nêutrons}, 17 \omterm{def:g9:inside-the-atom:nucleus}{elétrons}.

\textbf{3.} Os dois têm 17 \omterm{def:g9:inside-the-atom:proton}{prótons}, logo os dois são o elemento cloro;
eles só diferem pelo número de \omterm{def:g9:inside-the-atom:proton}{nêutrons}, logo são \omterm{def:g9:inside-the-atom:isotope}{isótopos}.

\textbf{4.} Sim: a química depende dos \omterm{def:g9:inside-the-atom:nucleus}{elétrons}, e os dois têm 17.

\textbf{5.} $35 \times 1.67 \times 10^{-27} = \qty{5.845e-26}{kg}$.

\textbf{6.} $17 \times 9.11 \times 10^{-31} = \qty{1.55e-29}{kg}$:
cerca de 4000 vezes menos que o \omterm{def:g9:inside-the-atom:nucleus}{núcleo}, desprezível.

\textbf{7.} $37 \times 1.67 \times 10^{-27} = \qty{6.179e-26}{kg}$.

\textbf{8.} $\qty{1}{g} = \qty{e-3}{kg}$; $10^{-3} / 5.845 \times
10^{-26} \approx 1.71 \times 10^{22}$ \omterm{def:g7:atoms-and-molecules:atom}{átomos}.

\textbf{9.} \qty{75.8}{\percent} e \qty{24.2}{\percent}.

\textbf{10.} $758 \times 5.845 \times 10^{-26} = \qty{4.430e-23}{kg}$;
$242 \times 6.179 \times 10^{-26} = \qty{1.495e-23}{kg}$.

\textbf{11.} $4.430 \times 10^{-23} + 1.495 \times 10^{-23} =
\qty{5.925e-23}{kg}$.

\textbf{12.} $5.925 \times 10^{-23} / 1000 \approx
\qty{5.93e-26}{kg}$: a massa média de um \omterm{def:g7:atoms-and-molecules:atom}{átomo} de cloro.
\end{solution}
